% !PRINTED 82
\weberpage{82}
%<<CONT
\gesperrt{holds then and only then, if at the same time the
congruences}
\weq{5}{x_1\equiv 0\ (\mathrm{mod}\ i_1),\quad x_2\equiv 0\ (\mathrm{mod}\ i_2),\quad \dots\ x_\nu\equiv 0\ (\mathrm{mod}\ i_\nu)}
\gesperrt{are satisfied.}

According to this, then, $J$ is the least positive number which, for all $b$, satisfies the congruence
\[
b^J\equiv 1\ (\mathrm{mod}\ m)
\]
or the least positive number for which, for all systems of the $\beta$,
\weq{6}{J\beta_1\equiv 0\ (\mathrm{mod}\ c_1),\quad J\beta_2\equiv 0\ (\mathrm{mod}\ c_2),\quad \dots\ J\beta_\mu\equiv 0\ (\mathrm{mod}\ c_\mu).}

From this there result the first conclusions about the composition of the number $m$.

a) \gesperrt{An odd prime $q$ which is not contained in $J$ can divide $m$ only simply.}

For if $m$ is divisible by $q^\varkappa$, then one of the quantities $c$ is, say, $c_1=\varphi(q^\varkappa)=q^{\varkappa-1}(q-1)$; hence if $\varkappa>1$, then $c_1$ is still divisible by $q$, and from (6) it then follows that all the $\beta_1$ would have to be divisible by $q$, which contradicts assumption IV.

b) \gesperrt{An odd prime $q$ which divides $J$ can be contained in $m$ at most once more than in $J$.}

For one can conclude in the same way that if $c_1=q^{\varkappa-1}(q-1)$, and $J$ is not divisible by $q^{\varkappa-1}$, all the $\beta_1$ would have to be divisible by $q$.

c) \gesperrt{If $J$ is odd, then $m$ must also be odd.}

For if $m$ is divisible by a power of 2, and hence at least by 4, then the index modulus corresponding to the number 2 is even, and the corresponding $\beta$ would by (6) all have to be even, contrary to requirement IV.

d) \gesperrt{If $J$ is divisible by a power of 2, then $m$ can contain the factor 2 at most twice more often than $J$.}

For if $m$ is divisible by $2^\lambda$ and $\lambda>2$, then two of the index moduli are
\[
c_1=2,\quad c_2=2^{\lambda-2}.
\]
%CONT>>
